\subsection{Composition of systems of formal power series}

Let A be a commutative ring; we shall say that a system

\[
\mathbf {f} = (f _ {1}, \dots , f _ {p}) \in (\mathrm{A} [ [ \mathrm{X} _ {1}, \dots , \mathrm{X} _ {q} ] ]) ^ {p}\tag{8}
\]

of formal power series in the \(X_{j}\) ( \(1 \leqslant j \leqslant q\) ), with coefficients in A, is without constant term if this is true of all the \(f_{j}\) . For every system (8) of formal power series and every system

\[
\mathbf {g} = (g _ {1}, \dots , g _ {q}) \in (\mathrm{A} [ [ \mathrm{X} _ {1}, \dots , \mathrm{X} _ {r} ] ]) ^ {q}\tag{9}
\]

of \(q\) formal power series without constant term, we shall denote by \(\mathbf{f} \circ \mathbf{g}\) (or \(\mathbf{f}(\mathbf{g})\) ) the system of formal power series \(f_{j}(g_{1},\ldots ,g_{q})\) ( \(1\leqslant j\leqslant p\) ) in

\[
\left. \left(\mathrm{A} \left[ \left[ \mathrm{X} _ {1}, \dots , \mathrm{X} _ {r} \right] \right]\right) ^ {p} \right.
\]

(Algebra, Chapter IV, § 5, no. 5). If

\[
\mathbf {h} = (h _ {1}, \dots , h _ {r}) \in (\mathrm{A} [ [ \mathrm{X} _ {1}, \dots , \mathrm{X} _ {s} ] ]) ^ {r}
\]

is a third system without constant term, then

\[
(\mathbf {f} \circ \mathbf {g}) \circ \mathbf {h} = \mathbf {f} \circ (\mathbf {g} \circ \mathbf {h}).\tag{10}
\]

For every integer m,

\[
\left(\mathbf {f} ^ {(m)} \circ \mathbf {g} ^ {(m)}\right) \circ \mathbf {h} ^ {(m)} = \mathbf {f} ^ {(m)} \circ \left(\mathbf {g} ^ {(m)} \circ \mathbf {h} ^ {(m)}\right)
\]

where \(\mathbf{f}^{(m)}\) , \(\mathbf{g}^{(m)}\) , \(\mathbf{h}^{(m)}\) denote the systems of polynomials consisting of terms of total degree \(\leqslant m\) in the systems of formal power series f, g, h. But clearly the terms of total degree \(\leqslant m\) in the series of \((\mathbf{f} \circ \mathbf{g}) \circ \mathbf{h}\) (resp. \(\mathbf{f} \circ (\mathbf{g} \circ \mathbf{h})\) ) are the same as in \((\mathbf{f}^{(m)} \circ \mathbf{g}^{(m)}) \circ \mathbf{h}^{(m)}\) (resp. \(\mathbf{f}^{(m)} \circ (\mathbf{g}^{(m)} \circ \mathbf{h}^{(m)})\) ), whence our assertion.

For every system (8), we shall denote by \(M_{\mathbf{f}}\) or \(M_{\mathbf{f}}(\mathbf{X})\) , the Jacobian matrix \((\partial f_i / \partial \mathbf{X}_j)\) ( \(1 \leqslant i \leqslant p, 1 \leqslant j \leqslant q\) ) where \(i\) is the index of the rows and \(j\) that of the columns; for two systems (8) and (9), where \(\mathbf{g}\) is without constant term,

\[
M _ {\mathbf {f} \circ \mathbf {g}} = (M _ {\mathbf {f}} (\mathbf {g})). M _ {\mathbf {g}}\tag{11}
\]

where \(M_{\mathbf{f}}(\mathbf{g})\) is the matrix whose elements are obtained by substituting \(g_j\) for \(\mathbf{X}_j\) ( \(1 \leqslant j \leqslant q\) ) in each series element of \(M_{\mathbf{f}}\) ; this formula is just a reformulation of formula (9) of Algebra, Chapter IV, § 5, no. 8. We shall denote by \(M_{\mathbf{f}}(0)\) the matrix of constant terms of the elements of \(M_{\mathbf{f}}\) ; then we deduce from (11) that

\[
M _ {\mathbf {f} \circ \mathbf {g}} (0) = M _ {\mathbf {f}} (0). M _ {\mathbf {g}} (0).\tag{12}
\]

Given an integer \(n > 0\) , we shall write

\[
\mathbf {1} _ {n} = \mathbf {X} = (\mathrm{X} _ {1}, \dots , \mathrm{X} _ {n}) \in (\mathrm{A} [ [ \mathrm{X} _ {1}, \dots , \mathrm{X} _ {n} ] ]) ^ {n},\tag{13}
\]

which will be considered as a matrix with a single column.

For every system \(\mathbf{f} = (f_1, \ldots, f_n) \in (\mathrm{A}[[\mathrm{X}_1, \ldots, \mathrm{X}_n]])^n\) , \(M_{\mathbf{f}}\) is a square matrix of order \(n\) ; we shall denote by \(\mathrm{J}_{\mathbf{f}}\) or \(\mathrm{J}_{\mathbf{f}}(\mathbf{X})\) its determinant and by \(\mathrm{J}_{\mathbf{f}}(0)\) the constant term of \(\mathrm{J}_{\mathbf{f}}\) , equal to \(\det(M_{\mathbf{f}}(0))\) ; if \(\mathbf{g} = (g_1, \ldots, g_n)\) is a system without constant term in \((\mathrm{A}[[\mathrm{X}_1, \ldots, \mathrm{X}_n]])^n\) , then, by (11) and (12),

\[
J _ {f \circ g} = J _ {f} (g) \cdot J _ {g}\tag{14}
\]

\[
\mathrm{J} _ {\mathbf {f} \circ \mathbf {g}} (0) = \mathrm{J} _ {\mathbf {f}} (0) \mathrm{J} _ {\mathbf {g}} (0).\tag{15}
\]

PROPOSITION 5. Let A be a commutative ring and \(\mathbf{f} = (f_1, \ldots, f_n)\) a system without constant term of \(n\) series in \(\mathrm{A}[[\mathrm{X}_1, \ldots, \mathrm{X}_n]]\) . Suppose that \(\mathrm{J}_{\mathbf{f}}(0)\) is invertible in A. Then there exists a system without constant term \(\mathbf{g} = (g_1, \ldots, g_n)\) of \(n\) series in \(\mathrm{A}[[\mathrm{X}_1, \ldots, \mathrm{X}_n]]\) such that

\[
\mathbf {f} \circ \mathbf {g} = \mathbf {1} _ {n}.\tag{16}
\]

This system is unique and

\[
\mathbf {g} \circ \mathbf {f} = \mathbf {1} _ {n}.\tag{17}
\]

The existence and uniqueness of \(\mathbf{g}\) follows from Algebra, Chapter IV, § 5, no. 9, Proposition 10, applied to the \(n\) formal power series

\[
f _ {i} (\mathrm{Y} _ {1}, \dots , \mathrm{Y} _ {n}) - \mathrm{X} _ {i} \quad (1 \leqslant i \leqslant n).
\]

It follows from (15) and (16) that \(\mathrm{J}_{\mathbf{f}}(0)\mathrm{J}_{\mathbf{g}}(0)=1\) and hence \(\mathrm{J}_{\mathbf{g}}(0)\) is also invertible. We conclude that there exists a system \(\mathbf{h}=(h_{1},\ldots,h_{n})\) of n series without constant term in \(\mathrm{A}[[\mathrm{X}_{1},\ldots,\mathrm{X}_{n}]]\) such that \(g\circ h=1_{n}\) ; from this relation and (16) it then follows, with the aid of (10), that

\[
\mathbf {h} = \mathbf {1} _ {n} \circ \mathbf {h} = (\mathbf {f} \circ \mathbf {g}) \circ \mathbf {h} = \mathbf {f} \circ (\mathbf {g} \circ \mathbf {h}) = \mathbf {f} \circ \mathbf {1} _ {n} = \mathbf {f}.
\]

Proposition 5 and formulae (10) and (15) show that under the law of composition \((\mathbf{f},\mathbf{g})\mapsto \mathbf{f}\circ \mathbf{g}\) the set of systems \(\mathbf{f} = (f_1,\dots ,f_n)\) of \(n\) series without constant term in \(\mathrm{A}[[\mathrm{X}_1,\dots ,\mathrm{X}_n]]\) for which \(\mathrm{J}_{\mathbf{f}}(0)\) is invertible in A, is a group.

